\section{Chapter Plan}

1. Introduction 
2. Econometric Model (Cherkonov paper) 
- 
3. Empirical application
4. Conclusion
5. Appendix
6. Biblography


Occupation/ Industry
HH formal
Occupation proportion


Econometric model
- The original method of decomposing differences in the mean of an outcome variable has been extended to distributional parameters other than the mean (Juhn, Murphy and Pierce (1993) and DiNArdo, Fortin and Lemieux(1996)). 


- As the true selection procedure is unknown, all selection correction models make some assumption about selection itself. 
-  (C. Machado, 2017) 
- In imputation method, selection on unobservables is assumed away based on observed covariates (Blau , Kahn, 2006, Johnson, Kitamura, 
Neal, 2000, Neal, 2004, Neal, Johnson, 1996; Olivetti,  Petrongolo, 2008)

- Second approach on selection argues that selection becomes negligible when  participation rates are high. Hence wage gaps can be estimated in smaller sample selected according to observed characteristics abeit different characteristics and employment choices suggest differences on unobserved observations as well.

The third approach of control function, which is known, as in parametric models (Heckman, 1974) or when unknown is estimated by semi-parametric or non parametric methods,  requires an exclusion restriction; an instrument Z that shifts employment but is unrelated to wages. 

The fourth approach of bounding uses restrictions motivated by economic theory to tighten the worst-case scenario on the wage gap. 

%insert related literature and refine to make the flow of four appraoch of selection better





